\section{An explicit conditional reduction for the arbitrary-size permanent endpoint}

\subsection*{Research gate QA033: a square residue tensor}
This section proves a reduction and records a tested uniform
candidate. It does not prove nonvanishing of the resulting
square determinant for arbitrary matrix size.
Work over \(\mathbb C\), let \(m\geq3\), and choose pairwise
distinct scalars \(a_1,\ldots,a_m\). Put
\[
 M_0(x,y)=xJ+y\operatorname{diag}(a_1,\ldots,a_m),\qquad
 t=x/y,\qquad \Delta(t)=\prod_{i=1}^m(a_i+t).
                                                               \tag{RJ.1}
\]
Here \(J\) is the all-ones matrix.
The transverse variable space will have dimension
\[
                  r_0=2m-1,\qquad
                  N=m(m-1).
                                                               \tag{RJ.2}
\]
Thus there are \(2m+1\) variables including \(x,y\).

\paragraph{The factorial transform and every first and second cofactor.}
On polynomials of any of the relevant bounded degrees, let
\[
                       \mathcal T(t^j)=j!\,t^j.
                                                               \tag{RJ.3}
\]
It is an invertible linear operator in characteristic zero.
For a set \(U\subseteq\{1,\ldots,m\}\), put
\(\Delta_U(t)=\prod_{i\notin U}(a_i+t)\).
The permanent and first cofactors of \(M_0(t,1)\) are
\[
 \begin{split}
 \operatorname{per}M_0&=\mathcal T\Delta,\\
 \partial_{ij}\operatorname{per}M_0
       &=\mathcal T\bigl(t^{|\{i,j\}|-1}\Delta_{\{i,j\}}\bigr).
 \end{split}                                                   \tag{RJ.4}
\]
For two entries \(e=(i,j)\), \(f=(k,l)\), their second
cofactor is zero if \(i=k\) or \(j=l\).
Otherwise, with \(U=\{i,j,k,l\}\), it is
\[
 \partial_e\partial_f\operatorname{per}M_0
             =\mathcal T\bigl(t^{|U|-2}\Delta_U\bigr).
                                                               \tag{RJ.5}
\]

\begin{proof}
After fixing \(s=0,1,2\) compatible entries in a permanent
term, \(m-s\) rows and columns remain. A diagonal entry
\(a_i\) can be selected precisely at an index outside
the union \(U\) of the deleted row and column indices.
Choose a subset \(D\) of those indices.
All remaining entries come from \(tJ\), and their
permanent is \((m-s-|D|)!\,t^{m-s-|D|}\).
Summing over \(D\) gives exactly
\(\mathcal T(t^{|U|-s}\Delta_U)\).
Repeated row or column choices in (RJ.5) are forbidden
by multilinearity. This proves every formula.
\end{proof}

\paragraph{An explicit first-jet kernel.}
For a matrix \(B=(b_{ij})\) of transverse linear forms,
apply \(\mathcal T^{-1}\) to the binary coefficients
of the suffix-linear term of \(\operatorname{per}(M_0+B)\).
By (RJ.4), dividing the result by \(\Delta\) gives
\[
           \sum_i\frac{b_{ii}}{a_i+t}
              +\sum_{i\ne j}
                   \frac{t\,b_{ij}}{(a_i+t)(a_j+t)}.
                                                               \tag{RJ.6}
\]
Its residue at \(t=-a_i\) is
\[
       L_i(B)=b_{ii}
                 -a_i\sum_{j\ne i}\frac{b_{ij}+b_{ji}}{a_j-a_i}.
                                                               \tag{RJ.7}
\]
The \(m\) residue coordinates identify the binary
degree-\((m-1)\) space with \(\mathbb C^m\).
Consequently the suffix-linear term vanishes precisely
when \(B\) takes its values in
\[
                         K=\bigcap_i\ker L_i.
                                                               \tag{RJ.8}
\]
The diagonal coefficients in (RJ.7) form an identity
matrix, so \(\dim K=m^2-m=N\).
Writing \(E_{ij}\) for a matrix unit, a basis is
\[
 K_{ij}=E_{ij}
       +\frac{a_i}{a_j-a_i}E_{ii}
       +\frac{a_j}{a_i-a_j}E_{jj},
                         \qquad i\ne j.
                                                               \tag{RJ.9}
\]
These basis vectors are indexed by directed edges.

\paragraph{The explicit residue bilinear forms.}
The residues of a polynomial of degree at most \(m-2\),
after division by \(\Delta\), lie in
\[
                H_m=\{(c_1,\ldots,c_m):\sum_i c_i=0\},
                                                               \tag{RJ.10}
\]
and give an isomorphism to this space. Indeed the sum
of residues is the coefficient of \(t^{-1}\) at infinity;
it is zero, and evaluation at all \(m\) distinct roots
proves injectivity.

For raw entries \(e=(a,b)\), \(f=(c,d)\), define
\[
 R_i(e,f)=
 \begin{cases}
  0,&a=c\text{ or }b=d\text{ or }i\notin U,\\[2pt]
  \displaystyle
       \frac{(-a_i)^{|U|-2}}
                 {\prod_{u\in U\setminus\{i\}}(a_u-a_i)},
       &\text{otherwise},
 \end{cases}
 \qquad U=\{a,b,c,d\}.
                                                               \tag{RJ.11}
\]
Extend this bilinearly to matrices, and set
\[
                     S_i(e,f)=R_i(K_e,K_f).
                                                               \tag{RJ.12}
\]
These are symmetric bilinear forms on \(K\);
by (RJ.5) and (RJ.10), \(\sum_iS_i=0\).
If \(B:V\to K\) is the transverse pencil, then the
transformed quadratic term has residues
\[
                         \frac12 B^{\mathsf T}S_iB.
                                                               \tag{RJ.13}
\]
Thus, up to harmless diagonal normalizations of symmetric
matrix coordinates, its differential is the square map
\[
 \begin{split}
 \mathcal D_B:\operatorname{Hom}(V,K)
       &\longrightarrow H_m\otimes\operatorname{Sym}^2V^*,\\
 C&\longmapsto
       \bigl(B^{\mathsf T}S_iC+C^{\mathsf T}S_iB\bigr)_i .
 \end{split}                                                   \tag{RJ.14}
\]
Both dimensions are
\[
 Nr_0=(m-1)\frac{r_0(r_0+1)}2=m(m-1)(2m-1).
                                                               \tag{RJ.15}
\]

\paragraph{A sufficient condition for the full endpoint.}
If \(\mathcal D_B\) is invertible for some \(a_i,B\),
the full permanent projection to transverse order two
in \(2m+1\) variables is dominant.

\begin{proof}
Varying the transverse coefficients within \(K\) changes
neither the binary term nor the suffix-linear term, by
(RJ.7)--(RJ.8), and (RJ.14) supplies all quadratic terms.
For arbitrary transverse variations, (RJ.7) supplies the
entire suffix-linear coefficient space. Their quadratic
components can be canceled by the preceding \(K\)-variations.
Finally the diagonal first cofactors
\(\mathcal T(\Delta/(a_i+t))\) form a basis in binary
degree \(m-1\). Varying their entries by \(x\) and \(y\)
therefore supplies all binary degree-\(m\) coefficients.
Cancel any resulting linear and quadratic components
using the variations already constructed.
This proves surjectivity of the full differential.
The polynomial coordinate-injection argument of M-013
then proves dominance. F-016 yields full permanent
multiplicity for every length at most \(2m+1\) and
third row at most two.
\end{proof}

No invertibility claim for all \(m\) is included in this
conditional statement.

\paragraph{Only pairs and triangles occur in the corrected tensor.}
If the unordered edges underlying \(e\) and \(f\) are
disjoint, then
\[
                              S_i(e,f)=0
                              \quad\hbox{for every }i.
                                                               \tag{RJ.16}
\]
If \(i\) is outside both edges, this follows immediately
from (RJ.11). Otherwise, relabel the disjoint edges as
\((i,j)\), \((k,l)\), and put \(t=-a_i\).
The contributions of \(E_{ij}\) and the diagonal
correction at \(i\) in (RJ.9) factor as
\[
 \left(\frac{t}{a_j-a_i}
                   +\frac{a_i}{a_j-a_i}\right)
 \left(
   \frac{t}{(a_k-a_i)(a_l-a_i)}
       +\frac{a_k}{(a_l-a_k)(a_k-a_i)}
       +\frac{a_l}{(a_k-a_l)(a_l-a_i)}
 \right)=0.
 \]
The diagonal correction at \(j\) contributes no residue
at \(i\). Reversing either edge does not change this
argument. This proves (RJ.16).

For a single edge \(\{i,j\}\), put \(A=a_i\), \(B=a_j\).
The nonzero block at residue \(i\), in the order
\((K_{ij},K_{ji})\), is
\[
 \frac1{(A-B)^3}
 \begin{pmatrix}
  2AB&-(A^2-4AB+B^2)\\
  -(A^2-4AB+B^2)&2AB
 \end{pmatrix}.
                                                               \tag{RJ.17}
\]
Residue \(j\) is its negative.
For three distinct vertices \(i,j,k\), put
\(A=a_i,B=a_j,C=a_k\) and
\(\Omega=(A-B)(A-C)(B-C)\).
The blocks needed at residue \(i\) are
\[
 \begin{split}
 S_i\bigl((K_{ij},K_{ji}),(K_{ik},K_{ki})\bigr)
       &=-\frac{A}{(A-B)(A-C)}
                        \begin{pmatrix}0&1\\1&0\end{pmatrix},\\
 S_i\bigl((K_{ij},K_{ji}),(K_{jk},K_{kj})\bigr)
       &=-\frac{A}{\Omega}
                        \begin{pmatrix}B&C\\C&B\end{pmatrix},\\
 S_i\bigl((K_{ik},K_{ki}),(K_{jk},K_{kj})\bigr)
       &=\frac{A}{\Omega}
                        \begin{pmatrix}B&C\\C&B\end{pmatrix}.
 \end{split}                                                   \tag{RJ.18}
\]
These follow by substituting the three-term vectors
(RJ.9) into (RJ.11), and specify the entire tensor
together with symmetry and relabeling.

The dual kernel of (RJ.14) is therefore completely
explicit. Represent a functional by symmetric bilinear forms
\(Q_1,\ldots,Q_m\) on \(V^*\), equivalently by elements
of \(\operatorname{Sym}^2V\), modulo adding the same
form to all of them. It is in the kernel exactly when
\[
                \sum_{i=1}^m
                    Q_i\bigl((S_iB)_e,\,\eta\bigr)=0
           \quad\hbox{for every directed edge }e
                         \hbox{ and every }\eta\in V^*.
                                                               \tag{RJ.19}
\]
One may set \(Q_m=0\) to remove the common-form ambiguity.

\paragraph{A uniform candidate and the current finite evidence.}
Use coordinates
\[
             u_1,\ldots,u_m,\quad v_1,\ldots,v_{m-1},
                         \qquad v_m=0.
\]
For \(i<j\), the tested off-diagonal labels are
\[
                   b_{ij}=u_i+2u_j,\qquad
                   b_{ji}=v_i+2v_j .
                                                               \tag{RJ.20}
\]
Complete the diagonal using (RJ.7), so \(B\) takes
values in \(K\). With \(a_i=i\), the square maps
in sizes \(m=3,4,5,6\) have full ranks
\[
                            30,\ 84,\ 180,\ 330.
 \]
The simpler lower labels \(b_{ji}=v_i\) also give full
square rank for \(m=3,4,5\), as recorded separately in
\path{agents/jet_review/check_residue_simplified_embedding.py}.
Those labels preserve all old entries when a last
vertex is added and are another candidate for induction.

The exact modular reproducer is
\path{agents/jet_review/structured_residue_tensor.py}.
It implements (RJ.9)--(RJ.14) and checks
\(\sum_iS_i=0\).
The tests concern this specific conditional reduction.
They are not evidence sufficient to promote (RJ.20)
to an arbitrary-size theorem.

\paragraph{The zero-parameter boundary and a limitation of first-order lifting.}
Set \(a_m=0\), keeping the other parameters nonzero
and pairwise distinct.
Then \(K=K_{\mathrm{old}}\oplus K_{\mathrm{star}}\),
where the first summand is spanned by edges avoiding
\(m\) and the second by edges incident to \(m\).
The last residue is a nondegenerate hyperbolic form
on the \(2m-2\) star coordinates and zero on the old
coordinates:
\[
 S_m(K_{im},K_{mi})=1/a_i,\qquad
 S_m(K_{im},K_{im})=S_m(K_{mi},K_{mi})=0 .
                                                               \tag{RJ.21}
\]
All other entries of \(S_m\) vanish.
For (RJ.20), the star labels are
\(b_{im}=u_i+2u_m\), \(b_{mi}=v_i\).
Their common kernel in \(V\) is the line generated by
\[
        w:\quad u_m=1,\quad u_i=-2\ (i<m),\quad v_i=0.
                                                               \tag{RJ.22}
\]
Consequently \(Q_m=w^2\), with the other \(Q_i=0\),
is a universal nonzero cokernel vector of (RJ.14),
under the usual symmetric-tensor dual identification.
The tested square maps at this boundary have corank one,
but uniqueness of that cokernel for arbitrary \(m\)
remains to be proved.

Even assuming uniqueness, a first-order perturbation
argument in \(a_m\) cannot supply the missing direction.
Indeed, let \(C\) be in the right kernel at \(a_m=0\).
Pairing its last quadratic variation with \(w\) and
an arbitrary vector shows
\[
                         C_{\mathrm{star}}(w)=0,
\]
because \(B_{\mathrm{star}}\) is surjective and the
star form is nondegenerate.
Hence \(B(w),C(w)\in K_{\mathrm{old}}\).
But \(S_m(a_m)\) restricts to zero on
\(K_{\mathrm{old}}\times K_{\mathrm{old}}\) for
every \(a_m\): its endpoint unions never contain \(m\).
Thus the left/right first-derivative pairing is zero.
Any perturbation proof must analyze at least the
second-order Schur complement. No first-order
transversality assertion is used here.

\paragraph{Geometric dictionary and remaining gate.}
The cofactor calculation is an explicit integral
coefficient-family construction, with factorials and
differences \(a_i-a_j\) inverted where indicated.
Its monomial grading and polynomial coordinate maps
are the SGA~3, I, Propositions~4.7.2--4.7.3,
pp.~34--36, comodule constructions; the finite
coefficient families use the framework of FGA,
Bourbaki~221, following Proposition~3.8,
pp.~265--266, as in M-013.
The all-copy consequence, if the square map is proved
invertible, is the already-checked F-016
characteristic-zero Schur-support argument.
The unresolved object is now the explicit dual kernel
(RJ.19), or equivalently the general determinant of
(RJ.14). No arbitrary generic tuple of quadrics has
been substituted for this actual permanent tensor.
